\documentclass[11pt]{article}
\usepackage[a4paper,margin=27mm]{geometry}
\usepackage{amsmath,amssymb,amsthm,mathtools}
\usepackage[T1]{fontenc}
\usepackage{lmodern,microtype}
\usepackage[colorlinks=true,linkcolor=blue,citecolor=blue,urlcolor=blue]{hyperref}
\newtheorem{theorem}{Theorem}[section]
\newtheorem{lemma}[theorem]{Lemma}
\newtheorem{proposition}[theorem]{Proposition}
\theoremstyle{definition}
\newtheorem{definition}[theorem]{Definition}
\theoremstyle{remark}
\newtheorem{remark}[theorem]{Remark}
\newcommand{\R}{\mathbb R}
\newcommand{\dd}{\,d\mu}
\newcommand{\norm}[1]{\left\|#1\right\|}
\newcommand{\intM}{\int_M}
\DeclareMathOperator{\Ric}{Ric}
\DeclareMathOperator{\Vol}{Vol}
\title{Mass conservation for fractional porous-medium flow\\
on noncompact Riemannian manifolds}
\author{Kaiwen Jin\\ \small\texttt{yui.ui.siki@gmail.com}}
\date{\today}
\begin{document}
\maketitle
\begin{abstract}
We prove mass conservation for the spectral fractional porous-medium
equation on complete noncompact Riemannian manifolds with Ricci curvature
bounded below and a Euclidean-dimensional Faber--Krahn inequality.
The result holds for every nonnegative weak dual solution that is
spatially integrable at each finite time, for arbitrary nonnegative
integrable initial data, all exponents \(m>1\), and all fractional
orders \(0<s<1\). Thus the prescribed weak solution class cannot lose
mass at spatial infinity. The proof first obtains a mass upper bound
from positivity of the integrated Green potential. A positive energy
identity after heat regularization then yields global integrability
of the nonlinear power on every strictly positive time interval.
A conservative linear resolvent converts this integrability into
mass equality at every finite time.
\end{abstract}

\section{Introduction}

Let \(M\) be a complete connected noncompact Riemannian manifold.
We consider the fractional porous-medium equation
\begin{equation}\label{eq:pme}
 \partial_tu+(-\Delta_M)^s(u^m)=0,\qquad m>1,\quad 0<s<1,
\end{equation}
where the fractional power is defined spectrally.
For a nonnegative density, preservation of its integral expresses
the absence of mass loss at spatial infinity. On a noncompact
manifold this conclusion requires an argument even when the
underlying linear diffusion preserves mass: the weak formulation
permits only compactly supported tests, and the nonlinear power
need not initially be globally integrable.

Berchio, Bonforte, Grillo and Muratori developed a weak dual solution
theory for \eqref{eq:pme} under a Ricci lower bound and a uniform
Faber--Krahn inequality \cite{BBGM}. Their formulation evolves the
Green potential of \(u\), and accommodates a weighted \(L^1\)
continuity condition. Their existence and smoothing results provide
solutions obtained through nonlinear semigroup constructions and
monotone approximation. Conservation of mass is listed among the
open questions in Section~7 of that work. An estimate established
for the constructed solutions cannot simply be assigned to every
solution satisfying the weak dual identity.

For manifolds with nonnegative Ricci curvature and a finite fractional
Green kernel, von Criegern, Grillo and Monticelli
\cite[Corollary~8.6]{CGM} prove mass conservation for weak dual
solutions obtained by monotone approximation from solutions associated
with bounded integrable data. The present theorem treats any spatially
integrable weak dual solution in the stated class and allows a general
Ricci lower bound.

Here we prove the following assertion. Suppose \(M\) has dimension
\(N\ge2\), its Ricci curvature is bounded below, and its first
Dirichlet eigenvalue obeys
\[
 \lambda_1(D)\ge c\,\Vol(D)^{-2/N},\qquad c>0,
\]
for every relatively compact smooth domain \(D\).
Then every nonnegative weak dual solution of \eqref{eq:pme}, with
nonnegative initial data in \(L^1(M)\) and with \(u(t)\in L^1(M)\)
at each finite time, satisfies
\[
 \intM u(t)\dd=\intM u(0)\dd\qquad\text{for every }t>0.
\]
The exact weighted solution definition and the full statement are
given in Section~\ref{sec:setting}. In particular the conclusion
concerns all solutions in that class; it does not require identifying
them with a particular construction.

The dual evolution says that the Green potential decreases by the
time integral of \(u^m\). That time integral is initially controlled
only locally. We first show that a nonnegative Green potential of
an integrable signed function forces its integral to be nonnegative.
The proof uses integrable upper-level truncations of the potential
and a conservative linear resolvent. This gives a mass upper bound
for the given weak dual solution.

The remaining issue is global integrability of \(u^m\).
Heat regularization makes the potential energy finite, even when the
initial unregularized energy is infinite. Symmetry of the Green kernel
gives an exact identity involving only nonnegative terms and endpoint
values. Summing it over time partitions yields a dissipation
inequality without differentiating the solution or assuming energy
regularity. A Green-potential estimate and interpolation make the
energy bound uniform at positive times. Once \(u^m\) is globally
integrable on positive time intervals, the same resolvent argument
gives equality of endpoint masses.

The result addresses conservation in the stated spatially integrable
class; it does not assert uniqueness in the larger weighted solution
space. Section~\ref{sec:setting} records the linear estimates and
normalization. Section~\ref{sec:order} proves the potential-order
argument, and Section~\ref{sec:energy} proves the energy estimate and
mass conservation. All integrals use Riemannian volume \(\mu\).
Operator powers are spectral, and \(C\) denotes a positive constant
depending on the fixed geometric and equation parameters; its value
may change between formulas.

\section{Setting and linear estimates}\label{sec:setting}

Throughout, \(M\) has dimension \(N\ge2\) and satisfies
\begin{equation}\label{eq:geometry}
 \Ric\ge-(N-1)k\,g,\qquad k\ge0,\qquad
 \lambda_1(D)\ge c\,\mu(D)^{-2/N},\quad c>0,
\end{equation}
on every relatively compact smooth domain \(D\).
The first eigenvalue is understood variationally, as the infimum of
\(\int|\nabla\phi|^2/\int|\phi|^2\) over nonzero
\(\phi\in C_c^\infty(D)\).

The same inequality holds on every relatively compact open set
\(\Omega\). Indeed, for each \(\phi\in C_c^\infty(\Omega)\)
choose a smooth compactly supported function in \(\Omega\) equal
to one near \(\operatorname{supp}\phi\). A regular superlevel set
between zero and one is a relatively compact smooth domain containing
this support; its existence follows from Sard's theorem
\cite[Theorem~7.2]{Sard}. If it is disconnected, apply
\eqref{eq:geometry} to its components. Since each component has
volume at most \(\mu(\Omega)\), summing their Rayleigh inequalities
and taking the infimum over \(\phi\) gives the assertion with the
same \(c\) and \(N\). Thus the Faber--Krahn condition in
\cite[Assumption~2.1]{BBGM} follows from \eqref{eq:geometry}.

Write \(A=-\Delta_M\) for its nonnegative self-adjoint Friedrichs
realization, \(P_t=e^{-tA}\), and \(L=A^s\).
The heat kernel is symmetric, nonnegative, and satisfies
\begin{equation}\label{eq:gaussian}
 p_t(x,y)\le Ct^{-N/2}\exp[-d(x,y)^2/(Ct)],\qquad t>0;
\end{equation}
see \cite[(2.3)]{BBGM}. The heat flow is sub-Markov and contracts
\(L^1\) and \(L^\infty\). Define the normalized Green kernel and
its positive integral operator by
\begin{equation}\label{eq:green}
 G_s(x,y)=\frac1{\Gamma(s)}\int_0^\infty
 p_t(x,y)t^{s-1}\,dt,\qquad
 Gf(x)=\intM G_s(x,y)f(y)\dd(y).
\end{equation}
For nonnegative inputs the integral is initially allowed to be
infinite. Notice that \(N>2s\) in the entire parameter range.

For \(x_0\in M\), put
\[
 \norm{f}_{1,x_0,G_s}
 =\int_{B_1(x_0)}|f|\dd+
 \int_{M\setminus B_1(x_0)}|f(x)|G_s(x,x_0)\dd(x).
\]
\begin{definition}\label{def:wds}
A nonnegative measurable function \(u\) is a weak dual solution
with initial datum \(u_0\ge0\) if, for every \(T>0\) and \(x_0\in M\),
\[
 u\in C([0,T];L^1_{x_0,G_s}),\qquad
 u^m\in L^1((0,T);L^1_{\mathrm{loc}}(M)),\qquad u(0)=u_0,
\]
and
\begin{equation}\label{eq:wds}
 \int_0^T\!\intM
 \bigl[Gu\,\partial_t\psi-u^m\psi\bigr]\dd\,dt=0
\end{equation}
for every smooth \(\psi\) compactly supported in \((0,T)\times M\).
\end{definition}
This is the potential formulation of \cite[Definition~2.1]{BBGM},
with the normalized kernel \eqref{eq:green} and with only the compact
smooth tests needed here.

\begin{theorem}\label{thm:main}
Assume \eqref{eq:geometry}, \(m>1\), and \(0<s<1\).
Let \(u_0\ge0\) belong to \(L^1(M)\), and let \(u\) be any
nonnegative weak dual solution in Definition~\ref{def:wds}.
If \(u(t)\in L^1(M)\) for every finite \(t>0\), then
\[
 \intM u(t)\dd=\intM u_0\dd
 \qquad\text{for every finite }t>0 .
\]
\end{theorem}

We use two linear potential estimates from \cite{BBGM}. First, for
each fixed ball \(B_\sigma(x_0)\),
\begin{equation}\label{eq:weighted}
 \norm{Gf}_{L^1(B_\sigma(x_0))}
 \le C_{\sigma,x_0}\norm{f}_{1,x_0,G_s};
\end{equation}
this is the estimate in Remark~2.1, using Lemma~3.2 there.
It implies continuity of \(Gu\) in \(L^1_{\mathrm{loc}}\).
Second, Lemma~5.2 gives, for \(1<p<N/(2s)\),
\begin{equation}\label{eq:hls}
 \norm{Gf}_{pN/(N-2sp)}\le C_p\norm{f}_p .
\end{equation}
All source locators here refer to arXiv version~2 of \cite{BBGM}.
Its displayed heat potential (2.6) omits \(1/\Gamma(s)\).
If that printed kernel is denoted by \(\widehat G_s\), then
\(\widehat G_s=\Gamma(s)G_s\). The source weighted norm and the
one above are equivalent, with comparison factors
\(\min\{1,\Gamma(s)^{-1}\}\) and
\(\max\{1,\Gamma(s)^{-1}\}\). Dividing the potential estimates
by \(\Gamma(s)\) gives \eqref{eq:weighted}--\eqref{eq:hls}.
We justify the exact spectral representation separately.

\begin{lemma}\label{lem:linear}
The heat semigroup \(P_t\) and the fractional semigroup \(S_r=e^{-rL}\)
preserve integrals on \(L^1(M)\). For nonnegative \(f\in L^1(M)\),
\begin{equation}\label{eq:potential}
 Gf=\int_0^\infty S_rf\,dr
\end{equation}
almost everywhere. If \(p_w=N/(N-2s)\), then
\begin{equation}\label{eq:weak}
 \mu\{G|f|>\lambda\}\le
 C(\norm{f}_1/\lambda)^{p_w},\qquad \lambda>0,
\end{equation}
and, for every \(\varepsilon>0\),
\begin{equation}\label{eq:heatpotential}
 \norm{P_\varepsilon G|f|}_\infty
 \le C\varepsilon^{s-N/2}\norm{f}_1,\qquad
 P_\varepsilon Gf=GP_\varepsilon f\quad(f\ge0).
\end{equation}
\end{lemma}
\begin{proof}
Bianchi--Setti's cutoffs \cite[Corollary~2.3]{BS}, with their curvature
exponent zero, give \(0\le\chi_R\le1\) in \(C_c^\infty(M)\),
equal to one on \(B_R(o)\), and
\(\norm{A\chi_R}_\infty\le C/R\), \(R\ge1\).
Their curvature constant may be chosen with \(\kappa^2\ge k\);
their support parameter is fixed independently of \(R\).
Spectral calculus gives
\[
 P_t\chi_R-\chi_R=-\int_0^tP_\rho A\chi_R\,d\rho .
\]
For bounded compact \(f\), symmetry and contraction therefore imply
\[
 |\langle P_tf-f,\chi_R\rangle|
 \le t\norm{f}_1\norm{A\chi_R}_\infty\longrightarrow0 .
\]
Dominated convergence gives \(\int P_tf=\int f\).
Density and contraction extend this to all \(L^1\) inputs.
In particular \(P_t1=1\) in the kernel sense.
The subordination formula for \(0<s<1\) is
\[
 S_rf=\int_0^\infty P_\rho f\,d\nu_r(\rho),\qquad
 \int_0^\infty e^{-\lambda\rho}d\nu_r(\rho)=e^{-r\lambda^s};
\]
see the linear semigroup discussion in
\cite[proof of Proposition~5.1]{BBGM}.
The measures \(\nu_r\) have total mass one, by letting
\(\lambda\downarrow0\). Tonelli now proves conservation for \(S_r\).

For \(\lambda>0\), the two scalar integrals
\[
 \frac1{\Gamma(s)}\int_0^\infty e^{-t\lambda}t^{s-1}\,dt
 =\lambda^{-s}=\int_0^\infty e^{-r\lambda^s}\,dr
\]
agree. Spectral calculus \cite[Section~3.1]{Teschl} applies first
to quadratic pairings with bounded compact inputs \(h\).
These pairings are finite: splitting the heat integral at one gives
\[
 \norm{G|h|}_\infty\le
 \frac1{\Gamma(s)}
 \left(\frac{\norm{h}_\infty}{s}
       +\frac{C\norm{h}_1}{N/2-s}\right).
\]
There is no \(L^2\) zero mode, since \eqref{eq:gaussian} and symmetry
give \(\norm{P_th}_\infty\le Ct^{-N/4}\norm{h}_2\).
An invariant \(L^2\) function is therefore zero.
Polarization identifies mixed compact pairings. Increasing bounded
compact truncations of a nonnegative \(L^1\) input and applying
Tonelli gives \eqref{eq:potential}.

To prove \eqref{eq:weak}, split \eqref{eq:green} at a time \(T>0\).
The short part has \(L^1\) norm at most \(CT^s\norm{f}_1\),
and the long part has \(L^\infty\) norm at most
\(CT^{s-N/2}\norm{f}_1\). For \(f\ne0\) choose \(T\) so that
the latter bound is \(\lambda/2\), and use Chebyshev's inequality
on the short part. This gives the exponent \(p_w\) in \eqref{eq:weak}.
The zero input is immediate. In particular these potentials are finite
almost everywhere.
Finally, Tonelli and the heat composition law give commutation and
\[
 \norm{P_\varepsilon G|f|}_\infty
 \le C\norm{f}_1\int_0^\infty
       (\varepsilon+t)^{-N/2}t^{s-1}\,dt
 \le C'\varepsilon^{s-N/2}\norm{f}_1 .
\]
The last integral is finite because \(0<s<N/2\).
\end{proof}

\section{Potential order and endpoint evolution}\label{sec:order}

We first extract information about mass from an ordered potential.
The potential itself need not be integrable.

\begin{lemma}\label{lem:order}
Let \(h\in L^1(M)\) be real and suppose \(W=Gh\ge0\) almost
everywhere. Then \(\intM h\dd\ge0\).
If in addition \(W\in L^1(M)\), then \(\intM h\dd=0\).
\end{lemma}
\begin{proof}
Fix \(\alpha>0\) and define the positive resolvent and its
Markov normalization by
\[
 R_\alpha=\int_0^\infty e^{-\alpha r}S_r\,dr,\qquad
 T_\alpha=\alpha R_\alpha .
\]
Lemma~\ref{lem:linear} implies
\begin{equation}\label{eq:markov}
 T_\alpha1=1,\qquad
 \intM R_\alpha f\dd=\alpha^{-1}\intM f\dd,\qquad
 \intM T_\alpha f\dd=\intM f\dd\quad(f\in L^1).
\end{equation}
The extension to nonnegative measurable functions is by monotone
kernel integration.
For \(f\ge0\) in \(L^1\), \eqref{eq:potential} and Tonelli yield
\[
 T_\alpha Gf
 =\int_0^\infty(1-e^{-\alpha v})S_vf\,dv
 =Gf-R_\alpha f.
\]
The coefficient follows by setting \(v=r+t\) and integrating
\(\alpha e^{-\alpha r}\) over \(0<r<v\).
All these potentials are finite almost everywhere by
\eqref{eq:weak}. Applying the identity to the positive and
negative parts of \(h\), dominated by \(G|h|\), gives
\begin{equation}\label{eq:resolvent}
 W-T_\alpha W=R_\alpha h=:g\in L^1(M).
\end{equation}

Since \(0\le W\le G|h|\), estimate \eqref{eq:weak} and \(p_w>1\)
show that for each \(\ell>0\),
\[
 \intM(W-\ell)_+\dd
 =\int_\ell^\infty\mu\{W>\lambda\}\,d\lambda<\infty.
\]
Positivity and \(T_\alpha1=1\), applied to
\(W\le(W-\ell)_++\ell\), imply
\[
 (T_\alpha W-\ell)_+\le T_\alpha(W-\ell)_+ .
\]
Both truncated functions are integrable, so \eqref{eq:markov} gives
\[
 \intM\bigl[(W-\ell)_+-(T_\alpha W-\ell)_+\bigr]\dd\ge0.
\]
The absolute value of the integrand is bounded by \(|g|\).
Letting \(\ell\downarrow0\) in dominated convergence gives
\(\int g\ge0\), hence \(\int h=\alpha\int g\ge0\).
If \(W\in L^1\), direct integration of \eqref{eq:resolvent}
and \eqref{eq:markov} gives \(\int g=0\).
\end{proof}

\begin{proposition}\label{prop:massbound}
Every solution in Theorem~\ref{thm:main} satisfies, for all
\(0\le a<b<\infty\),
\begin{equation}\label{eq:endpoint}
 Gu(a)-Gu(b)=W_{a,b}:=\int_a^b u(t)^m\,dt\ge0
 \quad\hbox{in }L^1_{\mathrm{loc}}(M)
\end{equation}
and
\begin{equation}\label{eq:massbound}
 \norm{u(b)}_1\le\norm{u(a)}_1\le M_0:=\norm{u_0}_1 .
\end{equation}
Moreover \(u(t)\to u_0\) strongly in \(L^1(M)\) as \(t\downarrow0\).
\end{proposition}
\begin{proof}
The defining weighted continuity implies local \(L^1\) continuity
of \(u\), by covering compact sets with finitely many unit balls.
Estimate \eqref{eq:weighted} gives the same continuity for \(Gu\).
Using product tests in \eqref{eq:wds} shows that
\(\partial_tGu=-u^m\) distributionally.
For each compact smooth spatial test, the right side is integrable
up to \(0\) and \(T\). Its indefinite integral is continuous;
subtracting it leaves a continuous scalar function with zero
distributional derivative, hence a constant. Equality against all
spatial tests gives \eqref{eq:endpoint} for every endpoint, including
zero. In particular no time representative is discarded.

Lemma~\ref{lem:order} applies to \(h=u(a)-u(b)\) and
\eqref{eq:endpoint}, and proves \eqref{eq:massbound}.
For initial strong continuity, fix \(\delta>0\) and a compact \(K\)
with \(\int_{M\setminus K}u_0\dd<\delta\).
For small \(t>0\), local convergence gives
\(\int_K|u(t)-u_0|\dd<\delta\), while
\[
 \int_{M\setminus K}u(t)\dd
 \le M_0-\int_Ku(t)\dd
 \le\int_{M\setminus K}u_0\dd+\delta<2\delta .
\]
It follows that \(\norm{u(t)-u_0}_1<4\delta\).
\end{proof}

\section{Energy and conservation}\label{sec:energy}

If \(M_0=0\), Proposition~\ref{prop:massbound} already proves
Theorem~\ref{thm:main}. Hence assume \(M_0>0\).
We regularize with the ordinary heat semigroup, which commutes
with the spectral Green operator.

\begin{proposition}\label{prop:energy}
For \(\varepsilon>0\), define
\[
 f_\varepsilon(t)=P_\varepsilon u(t),\qquad
 V_\varepsilon(t)=Gf_\varepsilon(t)=P_\varepsilon Gu(t),
 \qquad E_\varepsilon(t)=\intM f_\varepsilon(t)V_\varepsilon(t)\dd .
\]
These energies are finite, and for every \(0\le a<b<\infty\),
\begin{equation}\label{eq:dissipation}
 E_\varepsilon(b)+2\int_a^b\!\intM f_\varepsilon^{m+1}\dd\,dt
 \le E_\varepsilon(a).
\end{equation}
For every \(t>0\) they also satisfy
\begin{equation}\label{eq:energybound}
 E_\varepsilon(t)\le
 C M_0^{\,\frac{N(m-1)+2s(m+1)}{N(m-1)+2s}}\,
 t^{-\frac{N-2s}{N(m-1)+2s}},
\end{equation}
where \(C\) is independent of \(\varepsilon\).
\end{proposition}
\begin{proof}
By \eqref{eq:heatpotential} and \eqref{eq:massbound},
\begin{equation}\label{eq:finite}
 0\le V_\varepsilon(t)\le C\varepsilon^{s-N/2}M_0,\qquad
 E_\varepsilon(t)\le C\varepsilon^{s-N/2}M_0^2 .
\end{equation}
Applying the positive heat kernel to \eqref{eq:endpoint} gives
\begin{equation}\label{eq:smoothevolution}
 V_\varepsilon(a)-V_\varepsilon(b)
 =\int_a^bP_\varepsilon(u(t)^m)\,dt\ge0.
\end{equation}
This positive integral is bounded by \(V_\varepsilon(a)\);
global integrability of \(u^m\) has not been used.

The mixed terms
\(\int f_\varepsilon(a)V_\varepsilon(b)\) and
\(\int f_\varepsilon(b)V_\varepsilon(a)\) are finite by
\eqref{eq:finite} and the mass bound. They are equal by
symmetry of \(G_s\). Expanding the product and canceling them,
then using positive Tonelli in \eqref{eq:smoothevolution}, gives
\begin{align}
 E_\varepsilon(a)-E_\varepsilon(b)
 &=\intM(V_\varepsilon(a)-V_\varepsilon(b))
       (f_\varepsilon(a)+f_\varepsilon(b))\dd \notag\\
 &=\int_a^b\!\intM P_\varepsilon(u(t)^m)
       (f_\varepsilon(a)+f_\varepsilon(b))\dd\,dt .       \label{eq:positiveenergy}
\end{align}
In particular \(E_\varepsilon\) is nonincreasing.

For fixed \(\varepsilon,x\), the function
\(t\mapsto f_\varepsilon(t,x)\) is continuous.
On a fixed ball this follows from local \(L^1\) continuity
and boundedness of \(p_\varepsilon(x,\cdot)\).
The contribution of its complement is uniformly small, since
\eqref{eq:gaussian} bounds it by
\(2CM_0\varepsilon^{-N/2}e^{-R^2/(C\varepsilon)}\) outside
\(B_R(x)\).
Sum \eqref{eq:positiveenergy} over partitions of \([a,b]\)
whose mesh tends to zero. The endpoint multipliers converge
pointwise to \(2f_\varepsilon(t,x)\); all integrands are
nonnegative. Fatou gives
\[
 E_\varepsilon(b)+2\int_a^b\!\intM
 P_\varepsilon(u^m)f_\varepsilon\dd\,dt\le E_\varepsilon(a).
\]
Conservativeness makes the heat kernel a probability measure, so
Hölder gives \(P_\varepsilon(u^m)\ge(P_\varepsilon u)^m\),
also for unbounded inputs by monotone truncation.
If \(f_\varepsilon(t,x)=0\), then \(u(t)=0\) almost everywhere
for that kernel measure, so \(P_\varepsilon(u(t)^m)(x)=0\)
as well. Thus the product is unambiguous and
\eqref{eq:dissipation} follows.

Set \(q=2N/(N+2s)\) and \(q'=2N/(N-2s)\).
These exponents are conjugate and \(1<q<N/(2s)\).
Since \(f_\varepsilon\in L^1\cap L^\infty\), the estimate
\eqref{eq:hls} and Hölder give
\[
 E_\varepsilon\le C\norm{f_\varepsilon}_q^2
 \le C M_0^\gamma
       \left(\intM f_\varepsilon^{m+1}\dd\right)^\beta ,
\]
where
\[
 \eta=\frac{(m+1)(N-2s)}{2Nm},\qquad
 \beta=\frac{2\eta}{m+1}=\frac{N-2s}{Nm}\in(0,1),
 \qquad\gamma=2(1-\eta)>0 .
\]
Indeed \(1/q=1-\eta+\eta/(m+1)\), which is the
\(L^1\)--\(L^{m+1}\) interpolation identity.
Consequently \eqref{eq:dissipation} implies
\begin{equation}\label{eq:scalar}
 E_\varepsilon(b)+k_0\int_a^bE_\varepsilon(t)^r\,dt
 \le E_\varepsilon(a),\qquad
 r=1/\beta>1,\quad
 k_0=2C^{-1/\beta}M_0^{-\gamma/\beta}.
\end{equation}

Here a scalar comparison can be proved without an absolute
continuity assumption. For any finite nonnegative nonincreasing
\(E\) satisfying \eqref{eq:scalar}, monotonicity gives
\(k_0(b-a)E(b)^r\le E(a)\).
If \(E(t)>0\), iterate this inequality on
\([t/2^{j+1},t/2^j]\) to obtain
\[
 \log E(t)\le r^{-n}\log E(0)
 +\sum_{j=0}^{n-1}r^{-(j+1)}
       [(j+1)\log2-\log(k_0t)].
\]
Letting \(n\to\infty\) gives
\[
 E(t)\le2^{r/(r-1)^2}(k_0t)^{-1/(r-1)}.
\]
If \(E(t)=0\), or \(E(0)=0\), the same bound is immediate.
Apply it separately for each \(\varepsilon\), using
\eqref{eq:finite}. Since
\[
 \frac1{r-1}=\frac{N-2s}{N(m-1)+2s},\qquad
 \frac{\gamma}{1-\beta}
 =\frac{N(m-1)+2s(m+1)}{N(m-1)+2s},
\]
we obtain \eqref{eq:energybound}, uniformly in \(\varepsilon\).
\end{proof}

\begin{proof}[Proof of Theorem~\ref{thm:main}]
Fix any \(0<a<b<\infty\). Proposition~\ref{prop:energy} bounds
\(\int_a^b\intM f_\varepsilon^{m+1}\dd\,dt\) uniformly
in \(\varepsilon\). The heat semigroup is strongly continuous
on \(L^1(M)\). One may see this directly from its \(L^2\)
continuity on bounded compact inputs, local \(L^1\) convergence,
and conservativeness, using the compact tightness argument in
Proposition~\ref{prop:massbound}; density and contraction give
the general case.
Thus, for each \(t\), \(P_\varepsilon u(t)\to u(t)\) in \(L^1\).
The difference has norm at most \(2M_0\), so dominated convergence
in time yields convergence in \(L^1((a,b)\times M)\).
Take a subsequence with summable errors to obtain almost everywhere
convergence. Fatou and the uniform dissipation bound imply
\[
 \int_a^b\!\intM u^{m+1}\dd\,dt<\infty .
\]
Hölder in time and space, applied to
\(u^m=u^{1/m}(u^{m+1})^{(m-1)/m}\), now gives
\[
 \int_a^b\!\intM u^m\dd\,dt
 \le[(b-a)M_0]^{1/m}
 \left(\int_a^b\!\intM u^{m+1}\dd\,dt\right)^{(m-1)/m}
 <\infty .
\]
In particular the potential difference \(W_{a,b}\) in
\eqref{eq:endpoint} belongs to \(L^1(M)\).
The equality assertion of Lemma~\ref{lem:order}, with
\(h=u(a)-u(b)\), gives
\(\intM u(a)\dd=\intM u(b)\dd\).
All strictly positive-time masses therefore coincide.
Finally let \(a\downarrow0\), using the strong initial
\(L^1\) continuity in Proposition~\ref{prop:massbound}.
Their common value is \(\intM u_0\dd\).
Since \eqref{eq:endpoint} holds at every endpoint, this
proves the claimed equality for every finite \(t>0\).
\end{proof}

\paragraph{AI-use statement.}
OpenAI Codex was used to develop and check the arguments and to draft
and revise the text, including internal reviews in fresh AI contexts.
No machine-checked formalization or independent human verification
is claimed.

\begin{thebibliography}{9}
\bibitem{BBGM}
E. Berchio, M. Bonforte, G. Grillo and M. Muratori,
\emph{The fractional porous medium equation on noncompact
Riemannian manifolds}, Math. Ann. \textbf{389} (2024), 3603--3651.
\href{https://doi.org/10.1007/s00208-023-02731-6}{doi:10.1007/s00208-023-02731-6}.
Source version used for numbered locators:
\href{https://arxiv.org/abs/2109.10732v2}{arXiv:2109.10732v2} (2023).

\bibitem{CGM}
D.-E. von Criegern, G. Grillo and D. D. Monticelli,
\emph{Fractional porous medium equation on manifolds with nonnegative
Ricci curvature: existence of solutions and smoothing effects via
potential methods},
\href{https://arxiv.org/abs/2609.12954v1}{arXiv:2609.12954v1} (2026).

\bibitem{BS}
D. Bianchi and A. G. Setti,
\emph{Laplacian cut-offs, porous and fast diffusion on manifolds
and other applications},
\href{https://arxiv.org/abs/1607.06008v1}{arXiv:1607.06008v1} (2016).

\bibitem{Sard}
A. Sard,
\emph{The measure of the critical values of differentiable maps},
Bull. Amer. Math. Soc. \textbf{48} (1942), 883--890.
\href{https://webhomes.maths.ed.ac.uk/~v1ranick/papers/sard.pdf}{Original article}.

\bibitem{Teschl}
G. Teschl,
\emph{Mathematical Methods in Quantum Mechanics:
With Applications to Schr\"odinger Operators},
Graduate Studies in Mathematics \textbf{99},
American Mathematical Society, 2009.
\href{https://www.mat.univie.ac.at/~gerald/ftp/book-schroe/schroe.pdf}{Author-hosted 2009 edition}.
\end{thebibliography}
\end{document}
